% Source: https://artofproblemsolving.com/community/c3h3133907_what_i_learned_in_middle_school
% Stable identity: aops-topic-3133907-q03
A standing player and the players currently sitting because of that player form a cluster. The cluster sizes partition five. A tag from a cluster of size $a$ to one of size $b$ replaces them by a cluster of size $a+1$ and $b-1$ singleton clusters. Thus cluster sizes determine transition probabilities. Let $u,v,w,t,s$ denote expected remaining times from partitions $(2,1,1,1),(2,2,1),(3,1,1),(3,2),(4,1)$. Choosing ordered standing roots gives
\[u=1+\dfrac{1}{4}w+\dfrac{1}{4}u+\dfrac{1}{2}v,\quad v=1+\dfrac{1}{3}w+\dfrac{1}{3}t+\dfrac{1}{3}v,\]
\[w=1+\dfrac{1}{3}s+\dfrac{1}{3}u+\dfrac{1}{3}t,\quad t=1+\dfrac{1}{2}s+\dfrac{1}{2}w,\quad s=1+\dfrac{1}{2}u.\]
Substitution or elimination gives $(u,v,w,t,s)=(14,13,12,11,8)$. Starting with five singletons, the first tag always reaches $(2,1,1,1)$, so the answer is $\boxed{1+14=15}$. The expectations are finite: from any state, repeatedly choosing one fixed standing player as tagger removes another standing player each turn, reaching absorption in at most four turns with probability bounded below by $20^{-4}$. Consequently the tail is bounded by a geometric sequence, justifying the expectation equations.
